\documentclass[10pt]{article}
\usepackage[T1]{fontenc}
\usepackage[utf8]{inputenc}
\usepackage[french]{babel}
\usepackage{newpxtext,newpxmath}
\usepackage[a4paper,top=1.5cm,bottom=1.35cm,left=1.6cm,right=1.6cm]{geometry}
\usepackage{microtype}
\usepackage{xcolor}
\usepackage{titlesec}
\usepackage[most]{tcolorbox}
\usepackage{enumitem}
\usepackage{array}
\usepackage{booktabs}
\usepackage{fancyhdr}

\definecolor{navy}{HTML}{1F3A5F}
\definecolor{grey}{HTML}{54606E}
\definecolor{green}{HTML}{2C6E49}
\definecolor{amber}{HTML}{9C6A16}
\definecolor{red}{HTML}{A6322E}
\definecolor{plum}{HTML}{7A4EA8}
\definecolor{boxgray}{HTML}{F3F5F8}
\definecolor{boxred}{HTML}{FCF2F1}
\definecolor{boxgreen}{HTML}{F1F7F3}
\definecolor{boxplum}{HTML}{F7F3FB}
\definecolor{boxcode}{HTML}{F7F8FA}
\definecolor{boxedge}{HTML}{D9E0EA}

\setlength{\parindent}{0pt}
\setlength{\parskip}{2.6pt}
\linespread{1.0}

\pagestyle{fancy}\fancyhf{}
\lhead{\small\color{grey}Scella}
\rhead{\small\color{grey}Les ensembles -- internes à chacun des deux corps -- 11 août 2026}
\cfoot{\small\color{grey}\thepage}
\renewcommand{\headrulewidth}{0.3pt}

\titleformat{\section}{\color{navy}\large\bfseries}{}{0pt}{}
\titlespacing{\section}{0pt}{11pt}{3pt}
\titleformat{\subsection}{\color{navy}\bfseries}{}{0pt}{}
\titlespacing{\subsection}{0pt}{7pt}{1.5pt}

\newtcolorbox{these}{colback=boxgray,colframe=boxedge,boxrule=0.5pt,arc=2.5pt,
  left=8pt,right=8pt,top=4.5pt,bottom=4.5pt}
\newtcolorbox{alerte}{colback=boxred,colframe=red!55,boxrule=0.8pt,arc=2.5pt,
  left=8pt,right=8pt,top=4.5pt,bottom=4.5pt}
\newtcolorbox{acquis}{colback=boxgreen,colframe=green!45,boxrule=0.8pt,arc=2.5pt,
  left=8pt,right=8pt,top=4.5pt,bottom=4.5pt}
\newtcolorbox{piste}{colback=boxplum,colframe=plum!45,boxrule=0.8pt,arc=2.5pt,
  left=8pt,right=8pt,top=4.5pt,bottom=4.5pt}
\newtcolorbox{extrait}{colback=boxcode,colframe=boxedge,boxrule=0.5pt,arc=1.5pt,
  left=7pt,right=7pt,top=3.5pt,bottom=3.5pt}

\newcommand{\best}[1]{{\color{green}\textbf{#1}}}
\newcommand{\pire}[1]{{\color{red}\textbf{#1}}}
\newcommand{\sz}[1]{{\small\color{grey}\textsf{#1}}}
\newcommand{\code}[1]{\texttt{\small #1}}
\newcommand{\x}{$\times$}
\newcommand{\loc}[1]{{\footnotesize\color{grey}\texttt{#1}}}

\begin{document}

\begin{center}
{\Huge\bfseries\color{navy} Scella}\\[4pt]
{\large\color{grey}\itshape Les ensembles, internes à chacun des deux corps}\\[4pt]
{\small\color{grey} 11 août 2026 \; · \; suite de la série des couples \; · \; quatre lentilles intrinsèques, plus un garde de fermeture}
\end{center}
\vspace{2pt}

\begin{these}
\textbf{Le résultat central, et il contredit la question.} Vous demandez les meilleurs ensembles
\emph{internes} à Scella et à EFE. Ils existent et ils sont solides -- mais ils arrivent
\textbf{troisième et quatrième}. Les deux premiers sont des \textbf{greffes} : du scellement posé sur un
producteur. Chaque corps, seul, plafonne : Scella sans producteur est un principe, EFE sans scellement
est un calcul que personne ne peut contester. Et la mesure fait apparaître l'asymétrie qui explique
tout le reste : \textbf{Scella est le seul tissu qui traverse EFE ; EFE n'a aucun tissu propre.}
\end{these}

\section{1 · Pourquoi le niveau « ensemble » dit autre chose que le niveau « couple »}

Quatre études de couples précèdent celle-ci. Monter d'un cran n'a d'intérêt que si le cran produit une
observation qu'une paire ne peut pas porter, et c'est le cas pour une raison structurelle : un
\textbf{couple} donne une \emph{propriété}, un \textbf{ensemble} donne une \textbf{capacité qu'un tiers
exerce de bout en bout}. « Artefact scellé \x{} re-prouveur zéro-dépendance » énonce un motif ;
« produire, sceller, nommer, publier l'ancre, re-prouver, refuser la falsification » est une chose qui
tourne. Le second se démontre, le premier s'explique.

D'où une modification assumée du cadre. Les quatre lentilles intrinsèques du 5 août restent la
notation -- \textbf{maturité} rapportée à l'effort restant, \textbf{rareté} pure, \textbf{adversaire} en
revue diligente, \textbf{démontrabilité} devant un ingénieur sceptique. J'y ajoute non pas une
cinquième note mais un \textbf{garde binaire} : l'ensemble \emph{ferme}-t-il sa boucle, ou lui
manque-t-il une pièce absente ? Un ensemble ouvert ne se note pas utilement, il se décrit. Faire de la
fermeture une note l'aurait diluée dans une moyenne ; en faire un garde la laisse visible.

\begin{alerte}
\textbf{Faits et jugements, séparés -- comme dans toute la série.} La partition des corps, les
dépendances, les emplacements et les zéros sont des \textbf{faits}, lus dans le code du 11 août ;
l'annexe donne la commande qui rejoue chacun. Les \textbf{notes} sont un \textbf{jugement} structuré
pour être discutable. Un désaccord sur « rareté = 4 » est légitime et ne se tranche par aucune
commande ; un désaccord sur « \code{HopKind} : 0 occurrence » s'ouvre en une seconde.
\end{alerte}

\section{2 · La partition des deux corps, mesurée et non supposée}

La distinction que vous posez entre Scella et EFE existe dans le code. Elle se lit dans les noms de
crates, et surtout dans le graphe de dépendances.

\renewcommand{\arraystretch}{1.15}
\begin{center}\small
\begin{tabular}{@{}>{\raggedright\arraybackslash}p{2.9cm} c >{\raggedright\arraybackslash}p{11.8cm}@{}}
\toprule
\textbf{Corps} & \sz{Crates} & \textbf{Ce qu'il fait, et qui le compose} \\
\midrule
\textbf{Scella} \sz{le scellement} & 15 &
  \code{provenance-dag}, \code{provenance-capture}, \code{provenance-bridge}, \code{provenance-verify},
  \code{audit-api}, \code{audit-console}, \code{forensic-export}, \code{reg-evidence},
  \code{verify-mini}, \code{verify-meta-mini}, \code{efe-meta-verify}, \code{efe-audit-core},
  \code{efe-publisher-identity}, \code{efe-swarm-provenance}, \code{efe-llm-provenance} \\
\addlinespace[2pt]
\textbf{EFE} \sz{le calcul} & 19 &
  \code{efe-core}, \code{efe-kernels}, \code{efe-meta-swarm\{,-server,-proto\}},
  \code{efe-nano} + 5 crates nano, \code{efe-swarm-\{gateway,sdk,memory,topology,orchestrator,dashboard\}},
  \code{efe-gold-test-gate}, \code{efe-swarm-provenance} \\
\bottomrule
\end{tabular}
\end{center}

\code{efe-swarm-provenance} figure dans les \emph{deux} lignes, et \code{efe-llm-provenance} appartient
au scellement tout en ne servant que des chemins de calcul. Ce ne sont pas des erreurs de classement :
ce sont les \textbf{charnières}, et elles se chaînent --
\code{efe-llm-provenance} $\rightarrow$ \code{efe-swarm-provenance} $\rightarrow$
\code{provenance-bridge} $\rightarrow$ \code{provenance-dag}.

\begin{acquis}
\textbf{La mesure qui recadre tout.} \code{provenance-dag} compte \textbf{onze dépendants}, dont
\textbf{les deux binaires de calcul} -- \code{efe-swarm-gateway} et \code{efe-meta-swarm-server}. Et ils
ne se contentent pas de déclarer la dépendance : ils écrivent tous deux dans le \textbf{même type
d'ancre}, \code{provenance\_dag::AnchorIndex}, avec le même vocabulaire de racine.

Autrement dit : \textbf{Scella traverse EFE de part en part. EFE, lui, n'a aucun tissu qui lui soit
propre} -- aucun type, aucun chemin de données que ses deux binaires partageraient sans passer par le
scellement.
\end{acquis}

\section{3 · Le classement des dix ensembles}

Rangs de 1 à 10 par lentille ; \emph{Éc.} est l'amplitude entre le meilleur et le pire rang, c'est-à-dire
le signal de variance. La colonne \emph{Corps} applique le garde de fermeture : \emph{pur} = l'ensemble
se suffit à lui-même dans son corps ; \emph{greffe} = il enjambe les deux ; \emph{joignant} = il
relierait EFE à lui-même.

\renewcommand{\arraystretch}{1.16}
\begin{center}\small
\begin{tabular}{@{}r >{\raggedright\arraybackslash}p{6.75cm} l cccc c c@{}}
\toprule
& \textbf{L'ensemble} & \sz{Corps} & \sz{Mat} & \sz{Rar} & \sz{Adv} & \sz{Dém} & \textbf{Moy} & \sz{Éc.} \\
\midrule
\textbf{1} & \textbf{La boucle du reçu méta} \sz{Couche 3 servie + M2 + liaison v4 + ancre 3a + deux vérifieurs + démo} & \sz{greffe} & 2 & 4 & \best{1} & \best{1} & \textbf{2,00} & 3 \\
\textbf{2} & \textbf{La boucle pile B} \sz{S0$\rightarrow$S6 + producteur \code{bundle/:id} + ancre + \code{verify-mini}} & \sz{greffe} & \best{1} & 3 & 2 & 3 & \textbf{2,25} & 2 \\
\textbf{3} & \textbf{Le triangle de re-preuve croisée} \sz{un artefact, deux implémentations, \code{cross\_check}} & \sz{\textbf{Scella pur}} & 3 & 2 & 3 & 2 & \textbf{2,50} & \best{1} \\
\textbf{4} & \textbf{La contraction servie} \sz{$\Phi$ déterministe + canari Lipschitz + porte H1 + topologie fixe} & \sz{\textbf{EFE pur}} & 5 & \best{1} & 4 & 5 & \textbf{3,75} & \pire{4} \\
5 & Le Fil B servi \sz{reformulateur SLM + circuit 1-sonde + provenance par-hop + deadlines} & \sz{EFE pur} & 4 & 5 & 5 & 4 & \textbf{4,50} & \best{1} \\
6 & Le participant réel de bout en bout \sz{\code{efe-core} $\rightarrow$ gRPC $\rightarrow$ manifeste $\rightarrow$ DoD\#1} & \sz{EFE pur} & 7 & 8 & 6 & 6 & \textbf{6,75} & 2 \\
7 & La souveraineté d'egress \sz{boot guard + allowlist + PII + cap + reçu R4} & \sz{greffe} & 6 & 9 & 7 & 7 & \textbf{7,25} & 3 \\
8 & \textbf{La couverture universelle de re-preuve} \sz{les six chaînes orphelines --- latent} & \sz{Scella pur} & 8 & 7 & 8 & 8 & \textbf{7,75} & \best{1} \\
9 & \textbf{Fil C, la Colonne} \sz{nano + essaim + méta-essaim + Fil B, un seul reçu --- latent} & \sz{joignant} & 9 & 6 & 9 & 9 & \textbf{8,25} & 3 \\
10 & \textbf{Le runtime nano} \sz{six crates, aucun appelant --- \pire{cassé}} & \sz{EFE pur} & \pire{10} & \pire{10} & \pire{10} & \pire{10} & \textbf{10,00} & \best{0} \\
\bottomrule
\end{tabular}
\end{center}

\section{4 · Le résultat central : les meilleurs ensembles ne sont internes à personne}

C'est la réponse directe, et elle contredit la formulation de la question.

Les deux premiers -- la boucle du reçu méta et la boucle pile B -- sont des \textbf{greffes}. Retirez
EFE : Scella n'a plus rien à sceller, et une pile de provenance sans producteur est une bibliothèque.
Retirez Scella : EFE produit un vecteur de consensus que personne ne peut opposer à quiconque, et un
calcul non opposable n'est pas un actif.

Les meilleurs ensembles \textbf{purement internes} arrivent donc \textbf{troisième et quatrième}, et il
vaut la peine de dire lesquels précisément, parce que ce sont eux la réponse à « internes de chacun ».

\subsection{4.1 · Scella pur -- le triangle de re-preuve croisée \sz{(3\textsuperscript{e}, 2,50)}}

Un artefact scellé, \textbf{deux ré-implémentations indépendantes} du vérifieur -- l'une à zéro
dépendance, l'autre à soixante-dix -- qui doivent tomber d'accord, et le contrôle croisé qui épingle
leur accord. C'est le seul ensemble du dépôt qui ne dépend d'\textbf{aucune origine} : il vaut pour
n'importe quel artefact scellé, quel que soit ce qui l'a produit. C'est ce qui en fait le Scella le plus
pur du portefeuille.

Sa valeur profonde n'est pas cryptographique, elle est \emph{méthodologique} : deux implémentations
écrites séparément qui doivent s'accorder transforment un format en \textbf{spécification}. Une
divergence entre elles est un bug de spec révélé avant l'auditeur. Son écart de 1 dit ce qu'il est --
bon partout, exceptionnel nulle part, et c'est le profil d'une fondation.

\subsection{4.2 · EFE pur -- la contraction servie \sz{(4\textsuperscript{e}, 3,75)}}

$\Phi$ rendu déterministe, sa constante de Lipschitz \textbf{mesurée} en boîte noire par un estimateur
adverse à bisection, une porte de démarrage qui \textbf{refuse de servir} si la non-expansivité est
falsifiée, et une topologie fixe qui borne le tout. Il ne demande rien à Scella : c'est le seul
ensemble EFE de premier plan qui se ferme sans emprunter au scellement.

\section{5 · L'ensemble le plus haute-variance, et pourquoi il mérite un traitement à part}

La contraction servie porte un écart de \textbf{4} : première en rareté, cinquième en maturité comme en
démontrabilité. C'est, au niveau des ensembles, l'analogue du diagnostic porté sur la co-signature R3
dans la série des couples -- mais dans une version bien plus favorable, puisque celle-ci \emph{tourne}.

Ce que je ne sais pas retrouver ailleurs tient en une phrase : un moteur numérique qui \textbf{mesure}
sa propre garantie de contraction au lieu de la déclarer, et qui a \textbf{réfuté son propre
certificat} -- $\hat{L} = 0{,}69$ mesuré contre $0{,}17$ annoncé analytiquement, la borne mesurée
retenue et la borne théorique abandonnée. Puis un \emph{fail-open} trouvé en revue -- un \code{.ok()}
avalait les explosions de $\Phi$ en les traitant comme des hors-domaine -- et fermé par un état à trois
branches. Ce n'est pas une preuve, c'est un falsificateur ; la distinction est faite dans le code
lui-même.

Ce qui le plombe est aussi net. Il faut expliquer ce qu'est une constante de Lipschitz avant que
quiconque comprenne pourquoi la chose est remarquable. Et surtout : \textbf{son test de bout en bout ne
tourne pas dans le jeu de features par défaut} -- \code{m1d\_core\_serving\_e2e} rapporte zéro test
exécuté sans \code{core-integration}. La chose la plus rare du dépôt est celle qui se montre le plus
mal, et son silence par défaut est parfaitement cohérent avec le constat général : tout est opt-in.

\begin{piste}
\textbf{Son vecteur de démonstration n'est pas la théorie.} C'est la phrase \emph{« il a réfuté sa
propre note »}. Un ingénieur sceptique n'a pas besoin de comprendre Lipschitz pour comprendre qu'un
système qui contredit sa propre documentation et retient le résultat défavorable a été construit par
quelqu'un qui préfère la mesure au confort.
\end{piste}

\section{6 · Les ensembles joignants -- zéro ligne, mais pas terrain vierge}

C'est ici que le Sandwich trouve sa place, et la réponse est sans ambiguïté.

\renewcommand{\arraystretch}{1.15}
\begin{center}\small
\begin{tabular}{@{}>{\raggedright\arraybackslash}p{5.4cm} >{\raggedright\arraybackslash}p{4.2cm} >{\raggedright\arraybackslash}p{5.9cm}@{}}
\toprule
\textbf{Ensemble joignant} & \textbf{Ce qu'il ferait} & \textbf{État vérifié} \\
\midrule
\textbf{Sandwich C.1} & compresseur et vérifieur en hops scellés autour des appels frontière & \code{HopKind} : \pire{0 occurrence} \\
\textbf{Essaim médiateur} & chaîne déclarée de $N$ petits agents entre deux LLM & \code{SwarmMediator}, \code{EdgePlan} : \pire{0} \\
\textbf{Fil C, la Colonne} & nano + essaim + méta-essaim + Fil B sous un seul reçu & \pire{0 ligne} \\
\bottomrule
\end{tabular}
\end{center}

Les trois ensembles qui joindraient EFE \emph{à lui-même} n'existent pas. Le Fil B vit dans
\code{efe-swarm-gateway}, le méta-essaim dans \code{efe-meta-swarm-server} : deux binaires distincts,
aucun chemin de données partagé. La paire qu'on croit voir sur un schéma d'architecture n'est pas dans
le code.

\begin{acquis}
\textbf{Mais ils ne partiraient pas de zéro} -- et c'est ce que l'analyse par ensembles fait apparaître
là où celle par couples ne le pouvait pas. Les deux binaires partagent \emph{déjà} le substrat :
\code{provenance-dag}, \code{AnchorIndex}, le même vocabulaire de racine et d'ancre. Ce qui manque n'est
pas la fondation, c'est le \textbf{flux}. Le Fil C compte zéro ligne, mais son terrain est nivelé :
l'estimation « XL » héritée du 5 août mérite d'être revue à la baisse, exactement comme 3b l'a été
en découvrant que \code{forensic-export} portait déjà l'\code{AnchorProof}.
\end{acquis}

\section{7 · L'ensemble mort, et le gisement}

\textbf{Le runtime nano} ferme la marche avec un écart de \textbf{0} : dernier sur les quatre lentilles,
parfaitement cohérent, cohéremment mort. Six crates -- \code{efe-nano}, \code{-sdk},
\code{-app-coordinator}, \code{-preanalysis}, \code{-swarm-adapter}, \code{-inference-daemon} -- et le
daemon n'a \textbf{aucun appelant dans le dépôt}. C'est le seul ensemble dont je dirais qu'il coûte plus
qu'il ne rapporte tant qu'il n'est ni câblé ni retiré : du code jamais exécuté est une dette qu'un
examinateur trouve seul, et qu'il interprète comme un défaut de discipline plutôt que comme un
chantier.

\textbf{Le gisement} est au 8\textsuperscript{e} rang et vaut mieux que sa place. Neuf crates scellent
une véritable chaîne HMAC ; \textbf{deux seulement} ont un re-prouveur qu'un tiers peut lire
intégralement avant de lui accorder sa confiance. La propriété fondatrice de tout l'édifice couvre
\textbf{deux chaînes sur neuf}. Or le motif à répliquer est déjà écrit, et il est même classé
3\textsuperscript{e} : le triangle de re-preuve croisée existe deux fois, il pourrait exister neuf fois.
Aucune invention n'est requise -- seulement de la répétition disciplinée.

\begin{alerte}
\textbf{Une correction sur un ensemble que je classais comme acquis.} L'essaim $\Phi$ « pur en
concurrence » ne l'est \textbf{pas dans \code{HEAD}} : \code{dissolution\_centroid} y compte zéro
occurrence. Ces lignes vivent dans l'arbre de travail d'une session concurrente. Tout ensemble qui
s'appuie sur cette propriété est aujourd'hui \textbf{latent}, et non \emph{servi} -- la même correction
que celle portée par l'étude du croisement complet, re-vérifiée ici plutôt que recopiée.
\end{alerte}

\section{8 · Le jugement}

Vous demandiez les meilleurs ensembles internes à chacun des deux corps. Ils existent, ils sont
solides, et ils arrivent troisième et quatrième -- derrière deux greffes. Scella réalise sa valeur en se
posant sur un producteur ; EFE réalise la sienne en se faisant sceller. Aucun des deux ne vaut seul ce
qu'il vaut greffé, et la greffe, elle, tourne : \textbf{515 tests, zéro échec, zéro ignoré} sur la
chaîne du reçu, mesurés le 11 août sur 35 binaires de test.

Ce qui manque n'est donc pas un ensemble interne de plus. C'est qu'\textbf{EFE ne s'assemble jamais
avec lui-même}. Les trois ensembles qui le feraient -- Sandwich, essaim médiateur, Colonne -- sont à
zéro ligne, sur un terrain que le scellement a déjà préparé pour eux.

\begin{these}
\small\textbf{Rappel de cadrage, inchangé dans toute la série.} Le reçu atteste l'\textbf{intégrité}
d'une suite de pas, jamais la \textbf{correction} d'une décision. Il est \emph{tamper-evident pour un
détenteur de clé}, pas \emph{trustless pour un tiers}. Les participants du mode réel exécutent une
fusion latente Mode-1 numérique : ce ne sont pas des « vrais SLM ». Et ce document note des ensembles
sans présupposer aucune finalité extérieure -- il n'y en a plus de déclarée, et il n'en invente pas.
\end{these}

\section{Annexe · Vérifier chaque affirmation soi-même}

Les \emph{notes} de ce document ne se vérifient pas ; les \emph{faits} qui les portent, si. Comme dans
le reste de la série, les affirmations \textbf{négatives} sont les plus importantes à pouvoir contrôler :
une recherche qui revient vide n'est une preuve que si elle est rejouable.

\renewcommand{\arraystretch}{1.22}
\begin{center}\small
\begin{tabular}{@{}>{\raggedright\arraybackslash}p{5.3cm} >{\raggedright\arraybackslash}p{10.8cm}@{}}
\toprule
\textbf{Affirmation} & \textbf{Ce qui la rejoue} \\
\midrule
La partition des deux corps &
  \code{ls crates | grep -iE 'provenance|verify|audit|forensic|reg-evidence'} \newline
  \code{ls crates | grep -iE '\textasciicircum efe-(core|nano|swarm|meta-swarm|kernels|gold)'} \\
\code{provenance-dag} a onze dépendants &
  \code{grep -rl provenance-dag crates/*/Cargo.toml} \\
Les deux binaires écrivent la même ancre &
  \code{grep -rn 'provenance\_dag::AnchorIndex' crates/efe-swarm-gateway/src crates/efe-meta-swarm-server/src} \\
Les ensembles joignants sont à zéro &
  \code{grep -rc 'HopKind\textbackslash|SwarmMediator\textbackslash|EdgePlan' crates -{}-include='*.rs'} \\
Le daemon nano n'a aucun appelant &
  \code{grep -rl efe-nano-inference-daemon crates/*/Cargo.toml} \\
$\Phi$ pur n'est pas dans \code{HEAD} &
  \code{git show HEAD:crates/efe-core/src/els/mod.rs | grep -c dissolution\_centroid} \\
Le canari et la porte H1 existent &
  \code{crates/efe-meta-swarm/src/participants/lipschitz.rs} \newline
  \code{crates/efe-meta-swarm-server/src/participant/local.rs:342} \\
La topologie fixe existe &
  \code{crates/efe-core/src/genesis/fixed.rs:47} \\
Le Fil B est servi &
  \code{grep -rl EFE\_SLM\_REFORMULATOR\_URL crates -{}-include='*.rs'} \\
515 tests, zéro échec &
  \code{cargo test -p efe-meta-verify -p verify-meta-mini -p efe-meta-swarm-server} \\
\bottomrule
\end{tabular}
\end{center}

Deux avertissements, pour que l'annexe ne soit pas lue plus largement qu'elle ne vaut. Les lignes qui
portent sur des \emph{absences} valent à la date du document. Et la dernière dépasse dix minutes de
compilation au premier passage -- ce qui est précisément le genre de détail qu'il vaut mieux découvrir
avant une démonstration qu'au milieu.

\end{document}
